% Pursuing Stacks (cotes 134-2 to 134-8): archive cote 134-2, site page 53 (PDF page 54).
% appendix (letters to Breen). Draft compared: ps1.tex:1292-1354.
% AI-transcribed by Claude Opus 5.5 (claude-opus-5-5, adaptive thinking, effort high), Message Batches API, batch ps1 (msgbatch_01Rv5hCaNTtUt4JHx1yqoHPL), 2026-11-12; not yet checked by eye.
\scanpage{134-2/53}
\mspage{38} \mspage{6}

Mais on ne peut former \hand{$\mathcal{C}\mathcal{C}\mathcal{F} = \mathcal{C}^2\mathcal{F}$} et définir $H^{n+2}($ X , \hand{$\mathcal{F}$} ), semble-t-il
\uline{que} si \hand{$\mathcal{C}\mathcal{F}$} est lui- même un Gr-(n+1) champ, ce qui ne sera sans
doute le cas que si \hand{$\mathcal{F}$} est un n-champ de Picard strict ....

Venons en maintenant au cas où \hand{$\mathcal{F}$} est un n-champ \uline{localement}
\uline{constant} sur X, donc défini par un (n+1)- foncteur
\[ \hand{$\mathcal{C}$}_{n+1} \xrightarrow{\hand{$\mathcal{F}$}} \text{(n-Cat. de Picard strictes).} \]
Alors, posant pour $0 \hand{$\leq$} i \hand{$\leq$} n$
\[ H^i(\hand{$\mathcal{C}$}_{n+1}, \hand{$\mathcal{F}$}) \quad = \hand{$\pi$}_{n-i} (\underline{\text{Hom}} (e_n^{\hand{$\mathcal{C}_{n+1}$}}, \hand{$\mathcal{F}$})), \]
``on a fait ce qu'il fallait'' pour que l'on ait un isomorphisme canonique
\[ H^i(\hand{$\mathcal{C}$}_{n+1}, \hand{$\mathcal{F}$}) \hand{$\simeq$} H^i(\text{ X}, \hand{$\mathcal{F}$} ), \]
\hand{(}valable en fait sans structure de Picard sur \hand{$\mathcal{F}$} ...). Il s'impose, pour
tout \hand{$\infty$} -groupoide \hand{$\mathcal{C}$} et tout (n+1)- foncteur
\[ \hand{$\mathcal{C}$} \xrightarrow{\hand{$\mathcal{F}$}} \text{(n-cat. de Picard strictes)} \]
de définir les $H^i(\hand{$\mathcal{C}$}, \hand{$\mathcal{F}$})$, pour tout i, par
\[ H^i(\hand{$\mathcal{C}$},\hand{$\mathcal{F}$}) \quad = \hand{$\pi$}_{N-i} \underline{\text{Hom}} (e_N^{\hand{$\mathcal{C}$}} ,\hand{$\mathcal{C}$}^{N-n} \hand{$\mathcal{F}$}) \]
où on choisit N \hand{$\geq$}Sup (i,n). Si \hand{$\mathcal{F}$} n'a qu'une Gr-structure (pas
nécéssairement de Picard) on peut définir encore les $H^i($ \hand{$\mathcal{C}$}, \hand{$\mathcal{F}$} ) pour
i \hand{$\leq$} n+1 par
\margin{$\mathcal{C}\mathcal{F}$ (circled)}
\[ H^i(\hand{$\mathcal{C}$}, \hand{$\mathcal{F}$} ) \quad = \hand{$\pi$}_{n+1-i} \underline{\text{Hom}} ( e_{n+1}^{\hand{$\mathcal{C}$}} , \hand{$\mathcal{C}\mathcal{F}$}). \]
Dans le cas \hand{$\mathcal{C}$} = \hand{$\mathcal{C}$}$_{n+1}$ = \hand{$\Pi$}$_{n+1}$(X) , il doit être vrai encore
(en vertu de (A) plus haut), que cet ensemble est canoniquement isomorphe
à $H^{n+1}$(X, \hand{$\mathcal{F}$}) = \hand{$\pi$}$_o$ \hand{$\Gamma$}$_X$(\hand{$\mathcal{C}\mathcal{F}$}) (c'est vrai et bien facile pour n=0 ).
Décrire la flèche canonique entre les deux membres de
\[ H^{n+1}(X , \hand{$\mathcal{F}$}) \hand{$\simeq$} H^{n+1}(\hand{$\Pi$}_{n+1}X, \hand{$\mathcal{F}$}) \quad ? \]

Si on veut réexpliciter (A) et (B), en termes du yoga (C), on arrive
à la situation suivante:
On a un (n+1)- foncteur entre (n+1)-groupoides
\[ f_{n+1} : \hand{$\mathcal{C}$}_{n+1} \hand{$\longrightarrow$} \hand{$\mathcal{D}$}_{n+1} \]
induisant par troncature un n-foncteur

\mspage{53}
